%=====================================================================
\chapter{Reproducibility and Artefact Evaluation}\label{ch:reproducibility}
%=====================================================================
\locator{5}{Part~5 --- Rigour, Reproducibility and Ethics}

\begin{outcomes}
By the end of this chapter you will be able to:
\begin{tight}
  \item \textbf{Use} the ACM terms repeatability, reproducibility and
    replicability correctly.
  \item \textbf{Package} a study so a stranger can re-run it.
  \item \textbf{Capture} a computational environment rather than describing it.
  \item \textbf{Apply} the FAIR principles to your data and deposit it with a
    DOI.
  \item \textbf{License} code and data appropriately.
  \item \textbf{Proceed} responsibly when the data cannot be shared.
\end{tight}
\end{outcomes}

\begin{prereq}
\chref{mleval} for computational pipelines, \chref{benchmarking} for
non-determinism, \chref{qda} for the qualitative audit trail --- which is the
same idea in a different tradition.
\end{prereq}

\begin{keyterms}
repeatability $\bullet$ reproducibility $\bullet$ replicability $\bullet$
artefact $\bullet$ badging $\bullet$ container $\bullet$ environment pinning
$\bullet$ random seed $\bullet$ FAIR $\bullet$ persistent identifier $\bullet$
DOI $\bullet$ licence $\bullet$ synthetic data
\end{keyterms}

\section{Three Words That Are Not Synonyms}

ACM's definitions~\cite{acm2020badging} are now standard, and using them
loosely in a thesis is a small but noticeable error.

\begin{center}
\small
\begin{tabular}{@{}L{3.0cm}L{3.6cm}L{3.4cm}L{3.4cm}@{}}
\toprule
\textbf{Term} & \textbf{Who} & \textbf{Artefacts} & \textbf{Establishes} \\
\midrule
Repeatability   & The same team & The same setup & The measurement is stable \\
Reproducibility & A different team & The \emph{author's} artefacts &
                  The result follows from the data and code as given \\
Replicability   & A different team & \emph{Their own} artefacts &
                  The finding is about the world, not about the code \\
\bottomrule
\end{tabular}
\end{center}

\begin{conceptbox}[Reproducibility is a floor, not a goal]
A study that reproduces has established only that the reported numbers follow
from the deposited artefacts. That is a low bar --- and it is a bar a
substantial proportion of published computational work fails.

\medskip
Replication is the scientifically interesting property, and it is much harder:
new data, new implementation, and a finding that survives. Do not describe
having re-run your own scripts as ``replicating'' your result.
\end{conceptbox}

\begin{paperbox}[The crisis, and a programme for fixing it]
Marcus Munaf\`o and colleagues, \emph{A manifesto for reproducible science},
Nature Human Behaviour 1, 2017~\cite{munafo2017manifesto}.\oa{}

\medskip
Pair it with John Ioannidis, \emph{Why Most Published Research Findings Are
False}, PLoS Medicine 2(8), 2005~\cite{ioannidis2005why}\oa{}, and Monya Baker's
survey of 1,576 researchers~\cite{baker2016reproducibility}, in which more than
70\% reported having failed to reproduce another scientist's experiment and
more than half had failed to reproduce their own.

\medskip
\textbf{What to extract.} The manifesto's table of threats and measures ---
protecting against bias, improving methodological training, implementing
independent checks, encouraging transparency, diversifying peer review, and
changing incentives. Identify which of these you personally control, because
several of them you do.
\end{paperbox}

\section{What to Deposit}

\begin{center}
\small
\begin{longtable}{@{}L{3.0cm}L{4.8cm}L{6.4cm}@{}}
\toprule
\textbf{Component} & \textbf{Contents} & \textbf{Notes} \\
\midrule
\endfirsthead
\toprule
\textbf{Component} & \textbf{Contents} & \textbf{Notes} \\
\midrule
\endhead
\bottomrule
\endfoot
Data
  & Raw and processed, with a data dictionary
  & Both. Processed alone hides the cleaning; raw alone forces re-derivation \\
Code
  & Collection, cleaning, analysis, figure generation
  & Including the code that produced each figure and table, keyed to its number \\
Environment
  & Container image or lock file, with versions pinned
  & A prose description of ``Python 3.9 and scikit-learn'' is not an
    environment \\
Protocol
  & Pre-registration, instruments, consent forms
  & The design, not just the outcome \\
Seeds
  & Every random seed used
  & Without them nothing is repeatable, let alone reproducible \\
README
  & How to run it, expected runtime, expected output
  & Write it for a stranger. Then have a colleague follow it on a clean machine \\
Licence
  & For code and data separately
  & Undeclared means all rights reserved, which prevents reuse entirely \\
\end{longtable}
\end{center}

\begin{alertbox}[The one-command test]
Your artefact passes if a competent stranger, on a machine that is not yours,
can produce your headline result with one documented command and no
correspondence with you.

\medskip
Test it properly: a fresh container or a borrowed laptop, a colleague, and no
help from you. Almost every first attempt fails on an undocumented dependency,
an absolute path, or a data file that was never actually committed. Finding
that out now costs an afternoon; finding it out from an artefact evaluation
committee costs the badge.
\end{alertbox}

\section{Capturing an Environment}

\begin{center}
\small
\begin{tabular}{@{}L{2.4cm}L{4.0cm}L{7.7cm}@{}}
\toprule
\textbf{Level} & \textbf{Tool} & \textbf{Captures} \\
\midrule
Language deps & \texttt{requirements.txt} with pins, \texttt{renv.lock},
                \texttt{poetry.lock}
              & Library versions. Not the interpreter, the OS, or system
                libraries \\
Full stack    & Docker, Podman, Apptainer
              & OS, system libraries, interpreter, everything. The right answer
                for most computational work \\
Declarative   & Nix, Guix
              & Bit-reproducible builds. Strongest, steepest learning curve \\
Workflow      & Snakemake, Nextflow, \texttt{targets} in R
              & The dependency graph between steps, so partial re-runs are
                correct \\
\bottomrule
\end{tabular}
\end{center}

\begin{lstlisting}[language=,basicstyle=\ttfamily\footnotesize,caption={A Dockerfile
that will still build in three years. Every version is pinned, including the
base image digest --- tags move, digests do not.}]
FROM python@sha256:1e8a3f2c...   # pin the digest, not the tag
WORKDIR /work
COPY requirements.txt .
RUN pip install --no-cache-dir -r requirements.txt   # all versions pinned
COPY . .
# One command reproduces everything: analysis, tables, figures.
CMD ["bash", "run_all.sh"]
\end{lstlisting}

\begin{pitfall}[Seeds are necessary and not sufficient]
Setting a seed makes a single library's random draws repeatable. It does not
make results repeatable when:

\begin{tight}
  \item several libraries maintain independent random states --- Python's
    \texttt{random}, NumPy, and the framework each need seeding;
  \item GPU kernels use non-deterministic reduction orders, which most do by
    default;
  \item parallel execution merges results in completion order;
  \item hash randomisation changes dictionary iteration order between runs.
\end{tight}

\medskip
Seed everything you can, enable deterministic modes where the framework offers
them, and then \emph{report the variability that remains} across several seeds
rather than a single run's number. A result that only holds for one seed is not
a result (\chref{benchmarking}).
\end{pitfall}

\section{FAIR Data and Persistent Identifiers}

\begin{center}
\small
\begin{tabular}{@{}L{2.6cm}L{11.7cm}@{}}
\toprule
\textbf{F}indable    & A persistent identifier and rich metadata, indexed
                       somewhere searchable \\
\textbf{A}ccessible  & Retrievable by that identifier over an open protocol.
                       Metadata stays available even when the data itself is
                       restricted \\
\textbf{I}nteroperable & Standard formats and vocabularies; qualified
                       references to related data \\
\textbf{R}eusable    & A clear licence, provenance, and community standards
                       met~\cite{wilkinson2016fair} \\
\bottomrule
\end{tabular}
\end{center}

\begin{conceptbox}[Accessible does not mean open]
This is the most useful thing the FAIR principles say for anyone working with
human-subject or industrial data. A dataset can be FAIR and restricted: the
record, the metadata and the access procedure are public, while the data itself
requires an application.

\medskip
So ``my data is confidential'' does not exempt you. Deposit the metadata, the
instruments, the analysis code and a synthetic or aggregated version, and
document the access route. Then the study is checkable in every respect except
the one that genuinely cannot be.
\end{conceptbox}

\begin{center}
\small
\begin{tabular}{@{}L{2.4cm}L{4.2cm}L{7.5cm}@{}}
\toprule
\textbf{Repository} & \textbf{Best for} & \textbf{Notes} \\
\midrule
Zenodo   & Anything, up to 50 GB
         & CERN-operated, free, issues DOIs, integrates with GitHub so a
           release becomes a citable archive automatically \\
OSF      & Whole projects
         & Pre-registration, materials, data and preprints in one place, with
           version history \\
figshare & Figures, datasets & DOIs; some features commercial \\
Institutional & Long-term custody & Check what your university offers;
           it may be required for thesis data \\
\bottomrule
\end{tabular}
\end{center}

\begin{alertbox}[A GitHub link is not a deposit]
Repositories are deleted, renamed, made private and force-pushed. A URL in your
thesis that resolves today may not in two years, and the version it points at
may not be the one you used.

\medskip
Archive a release to Zenodo, cite the DOI, and put the GitHub URL beside it for
convenience. The DOI is the citation; the repository is the working copy.
\end{alertbox}

\section{Licensing}

\begin{center}
\small
\begin{tabular}{@{}L{2.2cm}L{4.4cm}L{7.5cm}@{}}
\toprule
\textbf{For} & \textbf{Common choice} & \textbf{Effect} \\
\midrule
Code & MIT, BSD-3, Apache-2.0
     & Permissive. Apache-2.0 adds an explicit patent grant \\
Code & GPL-3.0
     & Copyleft: derivatives must also be GPL \\
Data & CC0
     & Public domain dedication. Maximum reuse, no attribution requirement \\
Data & CC BY-4.0
     & Reuse with attribution. The common default \\
Text & CC BY-4.0
     & For the manuscript, where the publisher permits \\
\bottomrule
\end{tabular}
\end{center}

Do not apply a software licence to data or a Creative Commons licence to code:
each is drafted for the other kind of work and produces ambiguity. And note
that no licence at all is the most restrictive option, not the most permissive.

\section{Artefact Evaluation}

Most major computing venues now run artefact evaluation, awarding badges.

\begin{center}
\small
\begin{tabular}{@{}L{4.2cm}L{10.1cm}@{}}
\toprule
Artifacts Available   & Deposited in a public archive with a persistent
                        identifier. The easiest badge, and you should always
                        get it \\
Artifacts Evaluated --- Functional & Documented, consistent, complete, and it
                        runs \\
Artifacts Evaluated --- Reusable & Above, plus a standard of quality that lets
                        others build on it \\
Results Reproduced    & An independent committee obtained the paper's main
                        results using your artefacts \\
Results Replicated    & Obtained independently, without your artefacts \\
\bottomrule
\end{tabular}
\end{center}

\begin{templatebox}[C13 --- Reproducibility packaging checklist]
\begin{tabular}{@{}L{0.8cm}L{13.0cm}@{}}
$\square$ & Raw data deposited, or the access route documented if restricted \\
$\square$ & Processed data, with the script that produced it from raw \\
$\square$ & Data dictionary: every variable, its units, its coding \\
$\square$ & Analysis code, with each figure and table keyed to the code
            producing it \\
$\square$ & Environment captured as a container digest or lock file, all
            versions pinned \\
$\square$ & All random seeds recorded; variability across seeds reported \\
$\square$ & \texttt{run\_all} script producing every result with one command \\
$\square$ & README: requirements, runtime, expected output, hardware assumed \\
$\square$ & Licence for code; separate licence for data \\
$\square$ & Instruments, consent forms and protocol included \\
$\square$ & Pre-registration linked (\chref{prereg}) \\
$\square$ & Archived to Zenodo or OSF; DOI cited in the thesis \\
$\square$ & \textbf{Tested by a colleague on a clean machine, with no help from
            you} \\
\end{tabular}
\end{templatebox}

\begin{regbox}[Raw data is already required --- PhD \S14(i)]
``The raw data will have to be submitted to the supervisor along with a copy of
the thesis.''

\medskip
The regulation asks for the minimum; this chapter asks for the version that is
also usable. Having built a proper package for the supervisor, depositing it
publicly costs almost nothing more, and \reg{PhD Regs 2024}{\S14.3(B)} asks
whether ``the results of the study [are] reproducible'' as an assessment
criterion in its own right.
\end{regbox}

\begin{traditions}{reproducibility}
\tradEMP{Data and analysis scripts deposited; every analytic decision visible.
The bar is that another analyst gets your numbers from your data.}
\tradDSR{The artefact is the deposit, and badging is the natural venue
mechanism. Also deposit the evaluation harness --- without it the evaluation
cannot be checked.}
\tradQUAL{Raw transcripts usually cannot be shared, and the analogue is the
audit trail (\chref{qda}): codebook, memos, decision log, and de-identified
extracts.}
\tradFORM{Deposit proof scripts where machine-checked, and full proofs in an
appendix otherwise. A proof sketch that cannot be completed by a reader is not
reproducible either.}
\end{traditions}

\begin{lab}[Package and test]
\begin{tightnum}
  \item Work template C13 for a small analysis of your own.
  \item Write \texttt{run\_all.sh} that produces every figure from raw data.
  \item Containerise it, pinning the base image digest.
  \item Give a colleague only the archive and the README. Watch them run it and
    say nothing.
  \item Record every point at which they got stuck. Fix each, and repeat with a
    second colleague.
  \item Archive to Zenodo and put the DOI in your synopsis.
\end{tightnum}

\medskip
Step 4 is uncomfortable and is the whole exercise.
\end{lab}

\begin{checkpoint}
\begin{tightnum}
  \item Distinguish repeatability, reproducibility and replicability, and say
    which your own re-run of your scripts establishes.
  \item Why is setting a random seed insufficient for deterministic results?
  \item Your data contains personal information and cannot be shared. What can
    still be deposited, and what does FAIR require?
  \item Why is a GitHub URL an inadequate citation for an artefact?
\end{tightnum}
\end{checkpoint}

\begin{chaptersummary}
\begin{tight}
  \item Repeatability, reproducibility and replicability are distinct.
    Reproducibility is a floor; replication is the scientific goal.
  \item Deposit data, code, environment, protocol, seeds, README and licences.
    The one-command test is the standard, and it must be tested by someone else.
  \item Capture the environment as a container with pinned digests, not as
    prose.
  \item Seeds are necessary and insufficient. Report variability across seeds.
  \item FAIR data can be restricted: deposit metadata, instruments, code and a
    synthetic version, and document the access route.
  \item Archive to a repository with a DOI. A GitHub URL is a working copy, not
    a citation.
  \item Licence code and data separately. No licence is the most restrictive
    choice.
  \item \S14(i) already requires raw data for your supervisor, and \S14.3(B)
    assesses reproducibility.
\end{tight}
\end{chaptersummary}

\begin{reviewq}
  \item Artefact evaluation is voluntary at most venues and consumes reviewer
    effort. Should it be mandatory for empirical computing papers? Argue both
    sides, including the effect on scholars in resource-constrained
    institutions.
  \item Reproducibility guarantees that results follow from artefacts, not that
    they are true. Give a case where a fully reproducible result is
    nevertheless badly wrong, and say what would have caught it.
  \item Industrial collaboration frequently prohibits data sharing. Design a
    protocol allowing meaningful external verification without disclosing the
    data, and identify what it cannot verify.
  \item If the field required every paper to pass the one-command test, what
    would happen to publication volume, and would the trade be worth it?
\end{reviewq}
